\begin{tabularx}{\linewidth}{|>{\centering\arraybackslash}p{2.5cm}|>{\raggedright\arraybackslash}p{3cm}|>{\raggedright\arraybackslash}X|}

% --- CABECERA DE LA TABLA ---
% ltxtable la repetirá automáticamente en cada página.
\hline
\rowcolor{cyan!80}
\multicolumn{3}{|c|}{\textbf{Software Resources}} \\ \hline
\rowcolor{cyan!40}
\textbf{Resource ID} & \textbf{Resource Name} & \textbf{Description} \\ \hline
\endfirsthead % Esto es para que la cabecera se muestre en la primera página

% --- Definimos la cabecera para las páginas siguientes ---
\hline
\rowcolor{cyan!40}
\textbf{Resource ID} & \textbf{Resource Name} & \textbf{Description} \\ \hline
\endhead

% --- CONTENIDO DE LA TABLA ---
GIT & Git & Version control tool for managing the source code repository. \\ \hline
AZR & Microsoft Azure Resources & Cloud services and resources: SQL Database, Blob Storage, AI model, Event Grid, Service Bus, SignalR, EntraID, and Azure DevOps (for repository and product backlog management). \\ \hline
MCS & Microsoft Collaboration Tools & Communication and coordination platforms: Microsoft Teams, Calendar, Outlook email, Word editor. \\ \hline
VSC & Visual Studio Code & Source code editor with extensions and debugging tools. \\ \hline
FIB & FIB Platforms & Platforms from Barcelona School of Informatics(FIB): Racó, Atenea, FIB secretary, bachelor’s thesis official website guide. \\ \hline
GM & Gmail & Email communication with university entities. \\ \hline
WEB & Web Browsers and Search Engines & Research Browsers: Mozilla Firefox, Google Chrome, Microsoft Edge. \\ \hline
OVL & Overleaf LaTeX Editor & Thesis formatting editor for final thesis documentation. \\ \hline
OT & Online Tools & Diagram design tools: Edraw, Draw.io, and Gantt Project. \\ \hline

\end{tabularx}
